% ---
% filename : ["electrotechnique_circuits_dc_20260429_154_courant_circuit_mixte_resistances.tex"]
% id : "electrotechnique_circuits_dc_20260429_154"
% title : "Courant dans un circuit mixte de résistances"
% domain : "Électrotechnique"
% subdomain : "Circuits DC"
% tags : ["circuits dc", "résistance", "loi d'Ohm", "circuit mixte", "groupement parallèle", "groupement série", "ts"]
% difficulty : 2
% time_solve : 4  # Calcul : 2 (difficulte) * (1 question * 2)
% has_octave : true
% octave_files : ["electrotechnique_circuits_dc_20260429_154_courant_circuit_mixte_resistances.m"]
% author : "om"
% ---
\documentclass[a4paper,11pt,answers,addpoints]{exam}

% ---------------------------------------------------------
% LANGUE, ENCODAGE, FONTS
% ---------------------------------------------------------
\usepackage[utf8]{inputenc}

% ---------------------------------------------------------
% GÉOMÉTRIE ET MISE EN PAGE
% ---------------------------------------------------------
\usepackage[left=1.7cm,right=1.5cm,top=2cm,bottom=1cm]{geometry}
%\usepackage{a4wide}
\usepackage{microtype}      % Meilleure répartition du texte
\usepackage{float}          % Positionnement [H]
\usepackage{needspace}      % Saut conditionnel intelligent

% Éviter les veuves/orphelines (optionnel, à décommenter si besoin)
% \widowpenalty=10000
% \clubpenalty=10000

% Espacement vertical flexible
\raggedbottom
\setlength{\parskip}{0.5em plus 0.2em minus 0.1em}
\setlength{\parindent}{0pt}


% ---------------------------------------------------------
% MATHÉMATIQUES
% ---------------------------------------------------------
\usepackage{amsmath, amssymb, bm, esvect}

% Vecteurs
\newcommand{\V}{\overrightarrow}
\def\vect#1{\overrightarrow{\strut#1}}
\providecommand{\Degrees}[0]{\ensuremath{^\circ}}

% ---------------------------------------------------------
% UNITÉS ET GRANDEURS (SIunitx)
% ---------------------------------------------------------
\usepackage{siunitx}
\sisetup{output-decimal-marker={,}}
\DeclareSIUnit \var {var}
\DeclareSIUnit \VA {VA}
\DeclareSIUnit \kVA {kVA}
\DeclareSIUnit[quantity-product={}] \degree{\text{\textdegree}}

% ---------------------------------------------------------
% GRAPHIQUES : TikZ + CircuitikZ + PGFPlots
% ---------------------------------------------------------
\usepackage{tikz}
\usetikzlibrary{
	arrows.meta,
	intersections,
	decorations.pathreplacing,
	positioning
}

\usepackage[european,straightvoltages,EFvoltages]{circuitikz}

\usepackage{tkz-fct}
\usepackage{pgfplots}
\pgfplotsset{compat=1.12}

% Symbole étoile-triangle personnalisé
\newcommand{\wye}{%
	\mathbin{\tikz[x=1ex,y=1ex]{%
			\draw[line width=.1ex] (0,0)--(45:1)--++(-45:1) (45:1)--++(0,1);
	}}%
}


\usepackage{sankey}

\pgfdeclarelayer{background}
\pgfdeclarelayer{foreground}
\pgfdeclarelayer{sankeydebug}
\pgfsetlayers{background,main,foreground,sankeydebug}


% ---------------------------------------------------------
% TABLEAUX
% ---------------------------------------------------------
\usepackage{tabularx}
\usepackage{tabu}

% ---------------------------------------------------------
% HYPERLIENS
% ---------------------------------------------------------
\usepackage[colorlinks=true,
menucolor=red,
breaklinks=true,
urlcolor=blue,
linkcolor=blue]{hyperref}

% ---------------------------------------------------------
% COMMANDES PERSONNELLES
% ---------------------------------------------------------
\newcommand\nomauteur{bdminasmorgul@protonmail.com}
\newcommand\dateTe{29.04.2026}
\newcommand\brancheTe{TS}

% ---------------------------------------------------------
% STYLE DE L'EXERCICE
% ---------------------------------------------------------
\pagestyle{headandfoot}
\firstpageheadrule
\runningheadrule

\firstpageheader{\brancheTe\ (\nomauteur)}{Élève : }{(\dateTe) \thepage/\numpages}
\runningheader{\brancheTe\ (\nomauteur)}{Élève : }{(\dateTe) \thepage/\numpages}

% Compteur en bas de page
\newcommand{\totalpointtts}{%
	\begin{tikzpicture}[remember picture, overlay]
		\node[anchor=south west, inner sep=0pt, outer sep=0pt]
		at ([xshift=0.5cm,yshift=0.5cm]current page.south west)
		{\rotatebox{90}{Total des points : \numpoints}};
	\end{tikzpicture}
}

\firstpagefooter{\totalpointtts}{}{}
\runningfooter{\totalpointtts}{}{}

% Réglages exam
\printanswers
\addpoints
\pointsinmargin
\bracketedpoints

% ---------------------------------------------------------
% LISTINGS (code)
% ---------------------------------------------------------
\usepackage{xcolor}
\usepackage{listings}

\lstdefinestyle{octavestyle}{
	language=Matlab,
	basicstyle=\ttfamily\small,
	keywordstyle=\color{blue},
	commentstyle=\color{green!50!black},
	stringstyle=\color{red},
	stepnumber=1,
	frame=single,
	breaklines=true,
	breakatwhitespace=true,
	tabsize=4
}
% ---------------------------------------------------------
% DOCUMENT
% ---------------------------------------------------------
\begin{document}
	

	\unframedsolutions
	\pointsinmargin
	\bracketedpoints
\section*{Préparation TS PQ CFC ELMO, IELE, PELE}

\begin{questions}\question[9]
% Un micro-réseau de distribution industrielle est modélisé par le schéma de résistances ci-dessous. Les câbles et les récepteurs forment un circuit complexe. Calculer l'intensité du courant total $I$ circulant dans le circuit lorsqu'une tension de \qty{230}{[\volt]} est appliquée entre la borne d'entrée (point D) et la borne de sortie (point C)~?
	Dans une armoire de distribution, plusieurs barres de cuivre et traverses métalliques sont modélisées par un circuit de résistances.  
%	On applique une tension de $\qty{230}{[\volt]}$ entre la barre principale $D$ et la barre de retour $C$.
%		Le réseau de barres est modélisé par le schéma ci-dessous vu entre $D$ et $C$~:
		Calculer l'intensité du courant total $I$ circulant dans le circuit lorsqu'une tension de \qty{230}{[\volt]} est appliquée entre la borne d'entrée (point D) et la borne de sortie (point C)~?
\begin{center}
	\fbox{
		\begin{minipage}{0.85\linewidth}
			\textbf{Valeurs des résistances~:}
			\begin{align*}
				R_1 &= \qty{60}{[\ohm]}  & R_6 &= \qty{120}{[\ohm]} \\
				R_2 &= \qty{35}{[\ohm]}  & R_7 &= \qty{460}{[\ohm]} \\
				R_3 &= \qty{35}{[\ohm]}  & R_8 &= \qty{20}{[\ohm]}  \\
				R_4 &= \qty{80}{[\ohm]}  & R_9 &= \qty{1500}{[\ohm]} \\
				R_5 &= \qty{120}{[\ohm]} & R_{10} &= \qty{850}{[\ohm]}
			\end{align*}
			
		\end{minipage}
	}
\end{center}
\begin{center}
\begin{circuitikz}[scale=1,thick]
	% Entree D et premier segment
	%	\draw (0,0) node[left] {D} to[short, -*] (0.5,0);
	%	\draw (0.5,0) to[R, l=$R_7$] (2.5,0);
	
	% Branche superieure (R3 + R2 + R5)
	\draw (2.5,0) -- (2.5,1.5);
	\draw (2.5,1.5) to[R, l=$R_3$] (4.5,1.5);
	\draw (4.5,1.5) to[R, l=$R_2$] (6.5,1.5);
	\draw (6.5,1.5) to[R, l=$R_5$] (8.5,1.5);
	\draw (8.5,1.5) -- (8.5,0);
	
	% Branche inferieure (R6)
	\draw (2.5,0) -- (2.5,-1.5);
	\draw (2.5,-1.5) to[R, l=$R_6$] (8.5,-1.5);
	\draw (8.5,-1.5) -- (8.5,0);
	
	% Sortie R8 et R10
	\draw (8.5,0) to[R, l=$R_8$] (10.5,0);
	\draw (10.5,0) to[R, l=$R_{10}$,-o] (13.5,0) node[right] {C};

	\draw (0,0) node[left] {D} to[R,l=$R_7$, o-] (2.5,0);

	% Branches ouvertes
	\draw (4.5,1.5) to[R, l=$R_4$,-o] (4.5,4.5) node[above] {T};
	\draw (6.5,1.5) to[R, l=$R_1$,-o] (6.5,4.5) node[above] {A};
	\draw (10.75,0) to[R, l=$R_9$,-o] (10.75,-3) node[below] {B};
	
	\draw (0,-5) to[V, l_=$U_{DC}$, v={$230[V]$}] (13.5,-5);
	\draw (0,-5) to[short,-o] (0,0);
	\draw (13.5,-5) to[short,-o] (13.5,0);
\end{circuitikz}
\end{center}
\vspace{0.5em}
\needspace{55\baselineskip}
\begin{solution}
	L'analyse du chemin parcouru par le courant entre les points D et C permet d'identifier les résistances actives dans le circuit. Les branches aboutissant aux points A et B ainsi que la résistance $R_4$ sont en circuit ouvert par rapport au trajet D-C.
	
	\begin{align*}
		\text{Le trajet actif est le suivant : } &D \to R_7 \to (R_6 // (R_3 + R_2 + R_5)) \to R_8 \to R_{10} \to C \\[6pt]
		R_{branche\_gauche} &= R_3 + R_2 + R_5 = \qty{35}{[\ohm]} + \qty{35}{[\ohm]} + \qty{120}{[\ohm]} = \qty{190}{[\ohm]} \\[6pt]
		R_{parallel} &= \dfrac{1}{\dfrac{1}{R_{branche\_gauche}} + \dfrac{1}{R_6}} = \dfrac{1}{\dfrac{1}{\qty{190}{[\ohm]}} + \dfrac{1}{\qty{120}{[\ohm]}}} \\[6pt]
		R_{parallel} &= \qty{73,55}{[\ohm]} \\[6pt]
		\text{2. Calcul de la résistance équivalente totale } (R_{equ}) \text{ :} \\[6pt]
		R_{equ} &= R_7 + R_{parallel} + R_8 + R_{10} \\[6pt]
		R_{equ} &= \qty{460}{[\ohm]} + \qty{73,55}{[\ohm]} + \qty{20}{[\ohm]} + \qty{850}{[\ohm]} = \qty{1403,55}{[\ohm]} \\[6pt]
		\text{3. Calcul de l'intensité du courant } (I) \text{ :} \\[6pt]
		I &= \dfrac{U}{R_{equ}} = \dfrac{\qty{230}{[\volt]}}{\qty{1403,55}{[\ohm]}} \\[6pt]
		I &= \qty{0,16386}{[\ampere]} = \qty{163,86}{[\milli\ampere]} \\[6pt]
	\end{align*}
	\vspace{0.5em}
	\needspace{55\baselineskip}
	
	\begin{verbatim}
		% Télécharger OCTAVE depuis https://wiki.octave.org/Using_Octave et l'installer
		% Puis copier-coller le code dans OCTAVE
		
		% Données
		U = 230;            % Tension appliquee entre D et C [V]
		R7 = 460;           % Resistance segment entree [Ohm]
		R3 = 35;            % Resistance segment derivation 1 [Ohm]
		R2 = 35;            % Resistance segment derivation 2 [Ohm]
		R5 = 120;           % Resistance segment derivation 3 [Ohm]
		R6 = 120;           % Resistance charge en parallele [Ohm]
		R8 = 20;            % Resistance segment de liaison [Ohm]
		R10 = 850;          % Resistance segment sortie [Ohm]
		
		% Calcul de la branche en derivation
		R_br = R3 + R2 + R5;
		% Calcul du groupement mixte parallele
		R_par = 1 / (1/R_br + 1/R6);
		
		% Resistance totale du circuit de distribution
		Requ = R7 + R_par + R8 + R10;
		
		% Intensite du courant total par la loi d'Ohm
		I = U / Requ;
		
		% Affichage detaille des etapes de calcul
		printf("Resistance de la branche en derivation : %.2f [Ohm]\n", R_br);
		printf("Resistance equivalente du bloc parallele : %.2f [Ohm]\n", R_par);
		printf("Resistance totale equivalente (D-C) : %.2f [Ohm]\n", Requ);
		printf("Courant total circulant : %.4f [A] soit %.2f [mA]\n", I, I*1000);
	\end{verbatim}
	\end{solution}
\end{questions}
\end{document}
